\section{Introducción}
\label{sec:introduccion}

La conversación pública en redes sociales plantea un problema de clasificación
que no admite una lectura ingenua: el discurso de odio rara vez aparece
aislado. En un mismo tuit conviven la deshumanización de un grupo, la
agresión explícita y el señalamiento de una persona concreta. Los corpus
anotados para la tarea lo reflejan: en HatEval, por ejemplo, las etiquetas de
discurso de odio (\emph{hate speech}, HS), carácter dirigido (\emph{targeted},
TR) y agresividad (AG) se anotan de forma conjunta y condicionada
\cite{basile2019hateval}. Modelar las tres columnas como decisiones
independientes ignora esa estructura; modelarlas como una única clase de seis
combinaciones pierde la señal compartida y multiplica la escasez de cada
celda. La formulación multietiqueta, con un clasificador por etiqueta y una
regla ontológica que preserva la jerarquía, es el compromiso que adoptamos
aquí.

El contexto de este trabajo es la competencia académica \emph{TOXIC 2026 PUJ},
organizada para el curso de Procesamiento de Lenguaje Natural de la Maestría
en Inteligencia Artificial de la Pontificia Universidad Javeriana. La
competencia impone tres restricciones fundamentales: (i) prohibición de
utilizar arquitecturas basadas en transformers o modelos de lenguaje; (ii)
obligación de entrenar todas las representaciones desde cero a partir del
conjunto de entrenamiento; y (iii) prohibición de emplear el conjunto de
prueba para el ajuste del modelo base (la calibración de la capa de decisión
se guía únicamente por la retroalimentación agregada del leaderboard público). La métrica oficial es la media columnar de la $F_1$ binaria
de la clase positiva sobre HS, TR y AG \cite{opitz2019macro}. La nota final
combina el leaderboard privado con la claridad y justificación del cuaderno
entregado, de modo que el resultado y su explicación pesan por igual.

La pregunta que organiza este artículo es directa: \emph{¿puede un pipeline
clásico, sin transformers ni representaciones preentrenadas, competir en un
corpus pequeño y ruidoso de tuits en español?} La respuesta que documentamos
es afirmativa, pero con una distinción que consideramos central. El modelo
base ---representación subléxica, ensamble lineal y umbrales desacoplados---
alcanza un $F_1$ out-of-fold de 0.787582 y una mejora estadísticamente
sostenida sobre la línea base. La capa de decisión de última milla ---reglas
de staging y pares compensados--- eleva el leaderboard público hasta 0.81447,
pero lo hace con intervenciones específicas del test, guiadas por el propio
leaderboard. Separar ambas fuentes de mejora, en lugar de presentar el
0.81447 como si fuera una estimación limpia de generalización, es una de las
contribuciones metodológicas del trabajo.

Las contribuciones concretas son cinco. Primera, un pipeline clásico
completo y reproducible: una representación TF-IDF multiescala de 86\,275
características, un ensamble de tres optimizadores lineales con calibración
de temperatura y umbrales por etiqueta. Segunda, un protocolo de validación
agrupada que fusiona near-duplicados y pseudo-autores, verifica la ausencia
de fuga en cada pliegue y se replica en dos particiones congeladas. Tercera,
un análisis formal de la capa de decisión: demostramos que un par compensado
que rescata un verdadero positivo y suprime un falso positivo en la misma
columna incrementa la $F_1$ exactamente en $F_1/TP$, sin alterar el número de
positivos. Cuarta, una descomposición explícita de la ganancia del
leaderboard en sus dos componentes y una discusión de su validez externa.
Quinta, un contrato de reproducibilidad verificado: el envío campeón se
reconstruye bit a bit y dos ejecuciones independientes producen el mismo
archivo.

El resto del artículo sigue el orden de CRISP-DM. La
Sección~\ref{sec:trabajos} sitúa el trabajo frente a la literatura
(comprensión del negocio y estado del arte). La Sección~\ref{sec:datos}
describe el corpus y el análisis exploratorio (comprensión de los datos). La
Sección~\ref{sec:metodologia} detalla la preparación, la representación, los
modelos y la capa de decisión. La Sección~\ref{sec:resultados} recorre los
veinticinco experimentos evaluados en el leaderboard y separa lo que funcionó
de lo que no (modelado). La Sección~\ref{sec:discusion} presenta el análisis
estadístico, el análisis de error y los límites del enfoque (evaluación). La
Sección~\ref{sec:reproducibilidad} documenta el despliegue y la verificación
criptográfica. La Sección~\ref{sec:conclusiones} cierra con las lecciones y
el trabajo futuro.

Por convención numérica, y para no alterar los identificadores del registro
oficial de Kaggle, en todo el artículo se usa el punto como separador decimal
y el espacio fino para los miles.
